\documentclass{zhvt-classic}

% 使用本地源流明体（GenRyuMin2 TW）按文件路径加载，不依赖系统字体名
% 版框（内框／界栏／书口侧栏／鱼尾）改用与版心印章同一印泥棕色：
%   印章图片 yshs.png 实色 = RGB(135,68,52) = HTML 874434
\usepackage{xcolor}
\usepackage{tikz}
% sealbrown 双注册：xcolor 供 TikZ/pgf，L3 供 cls 版框
\definecolor{sealbrown}{HTML}{874434}
\ExplSyntaxOn
\color_set:nnn { sealbrown } { HTML } { 874434 }
\ExplSyntaxOff

\zhvtset{
  grid color       = sealbrown,
  main font        = GenRyuMin2TW-SB,
  main font path   = C:/Users/Administrator/texmf/fonts/opentype/public/GenRyuMin2TW/,
jiazhu font      = GenRyuMin2TW-M,
  jiazhu font path = C:/Users/Administrator/texmf/fonts/opentype/public/GenRyuMin2TW/,
  % 页码框容纳 4 字：正常页「九十七」+前置空 1 格=4 行恰好；
  % 100+ 页「一百零一」共 5 行超出，末字「一/二/三」被框裁掉（2026-10-04 实测），故放宽到 5
  shukou page number chars = 5,
  % 书口书名框高从默认 6 字减到 4 字，把省出的 2 字高拨给 chapter 框，
  % 使 6 字书名+卷名放进 9 字高的 chapter 框后不再压线（2026-10-05）
  shukou title chars = 4,
}

\title{史记}[日者列传]
\author{汉·司马迁}
\maker{闲叟漫记·重製}


% ===== 索隐述赞样式（2026-11-08 四次定稿 · 愚大人亲手改定参数）=====
% ⚠ 竖排前提：自定义字体家族必须带 RawFeature={vertical:+vert}；
%   simhei.ttf 加 +vhal 会整体加载失败静默回退，只写 +vert；simfang 同理。
\newCJKfontfamily{\shuzanfont}[Path=C:/Windows/Fonts/,RawFeature={vertical:+vert}]{simhei.ttf}
\newCJKfontfamily{\shuzanbody}[Path=C:/Windows/Fonts/,RawFeature={vertical:+vert}]{simfang.ttf}
% 题头：152x30pt 圆角 6pt 棕底白字；baseline (0,2) 使整框落入界行格子不压版框线；
% 白字节点 x 不补偿、y +8.2pt（像素级墨迹实测 dx/dy≈0，L/R 留白约 20pt）。
\newcommand{\shuzanhead}[1]{%
  \begin{tikzpicture}[baseline={(0,2)}]
    \fill[sealbrown, rounded corners=6pt]
      (-76pt,-15pt) rectangle (76pt,15pt);
    \node[text=white, inner sep=0pt, anchor=center] at (0pt,8.2pt)
      {\shuzanfont\fontsize{23}{23}\selectfont #1};
  \end{tikzpicture}%
}
% ★ 修正书口卷名位置：chapter 框起始位置改回 cls 默认值，避免全局影响。
%   只在每个卷首页通过 \juan 命令里的空命令微调该页位置（安全，不会跨页污染）。
\ExplSyntaxOn
\AtBeginDocument{
  \dim_set:Nn \l__zhvt_shukou_chapter_title_start_dim
    { \l__zhvt_shukou_title_height_dim + \l__zhvt_font_size_dim * 2 }
  \dim_set:Nn \l__zhvt_shukou_page_number_start_dim
    { \l__zhvt_shukou_chapter_title_start_dim + \l__zhvt_shukou_chapter_title_height_dim }
}
\ExplSyntaxOff

\begin{document}

% ★★ 版心（书口）排布覆写：把书名放到鱼尾之下（2026-10-04）
%
%   cls 默认版心自上而下：书名（title 静态框，在鱼尾之上）→ 鱼尾 → 篇名（chapter 动态框）→ 页码
%   本模板改为：          鱼尾（位置不动）→ 书名 → 篇名 → 页码分割线 → 页码
%
%   手法：把 cls 放在 title 静态框里的书名，改放进 chapter 动态框（鱼尾之下）；
%         title 静态框留空。书名取 \l__zhvt_title_str，即 \title{...} 的第一参数，
%         所以换书时只改 \title，不必在此重复写书名。
%
%   ⚠ 三条必须遵守的写法（2026-10-04 实测）：
%     ① 汉字必须用 \hbox{...} 包住。flowfram 会把框内容装进一个宽 14.4mm（= 一个 line_sep
%        = 书口栏宽）的纵向 minipage 并按该宽度折行，而整页靠 PDF /Rotate 90 显示、
%        框内文字本身不变向 —— 不包 \hbox 时「围炉夜话」4 字会被逐字折成 4 行，
%        在版面上就成了横排 4 字并横向溢出版框。
%     ② 竖排下**一个 \hbox 就是一整列**。书名与篇名要同处一列，必须放进同一个 \hbox；
%        分成两个 \hbox 时，后一个会被当成新的一列、甩到隔壁栏去
%        （实测「序」跑到书口栏外侧 14.4mm 的正文列上）。中间的 \quad 在框内是横向空白，
%        显示出来正好是书名与篇名之间的纵向间隔。
%     ③ ★ 内容盒要写成 **与框等宽的定宽盒** `\hbox to \l__zhvt_shukou_chapter_title_height_dim{…\hfil}`。
%        实测：chapter 框是把内容盒沿纵向**居中**摆放的，内容宽度一变，书名就跟着漂：
%        自然宽度的 \hbox（书名 38.4mm）落在 x=157.2mm，加篇名后（57.6mm）落在 x=152.4mm，
%        目录页（无篇名）与正文页的书名就差半格。做成等宽盒后，无论有没有篇名，
%        内容固定从框顶起排，实测两页均落在 x=142.8mm，不再漂。
\ExplSyntaxOn
% 把书口 chapter 框（书名+篇名）起始位置下移 1 字高，避免 4 字卷名顶到鱼尾
% ⚠ cls 的 \__zhvt_calculate_dimensions: 在 \AtBeginDocument 里执行，会重算该 dim；
%   必须把覆写挂到 begin document 之后，否则会被二次覆盖（2026-10-05 实测）。
% (旧方案已移到导言区，直接修改 chapter 框起始位置)
% 卷标题命令 #1=书口篇名(鱼尾) #2=正文卷名（含卷序，卷名内空格用 ~）
%
% ★★ 目录页两项修订（2026-10-05，愚大人指出）★★
%   ① 页码起点统一对齐：
%      原状 —— 目录项写成 \par\hspace{4em}卷名\hspace{4em}页码，页码起点 = 卷名实宽 + 4em。
%      卷名实宽随字数变：「卷一~梁惠王上」6 字 +~、「卷七~离娄上」5 字 +~，
%      差 1 个字位 → 卷七~卷十 4 条的页码整体比其余 10 条高 27.21pt（一个字），看着上下不齐。
%      改法 —— 卷名补全角空格（宽 = 1 字 = 27.215pt）补足到 **7 字符位**再写进目录，
%      14 条页码起点遂全部落在同一条竖线上（起点对齐；末字随页码字数自然参差）。
%      ⚠ 只对目录里的副本 \l_tmpc_tl 补，正文卷名仍用 \l_tmpa_tl 原样 —— 正文像素不受影响。
%   ② 卷名整体上移半字位：前导 \hspace 由 4em 改 3.5em。
%      实测原状目录块在版框内偏下（天头留白 139.6pt / 地脚留白 110.7pt），
%      上移 0.5em = 13.6pt 后上下基本均分，视觉居中。
\NewDocumentCommand { \juan } { mm } {
	\clearpage
	\tl_set:Nn \l__zhvt_chapter_title_tl { #1 }
	\tl_set:Nn \l_tmpa_tl { #2 }
	\tl_set_eq:NN \l_tmpc_tl \l_tmpa_tl
	\int_step_inline:nn { 7 - \str_count:n { #2 } }
	  { \tl_put_right:Nn \l_tmpc_tl { 　 } }
	\tl_set:Nx \l_tmpb_tl { \zhnum { page } }
	\tl_put_left:Nn \l_tmpb_tl { \hspace { 4em } }
	\tl_put_left:Nn \l_tmpb_tl { \l_tmpc_tl }
	\tl_put_left:Nn \l_tmpb_tl { \par \hspace { 3.5em } }
	\addtocontents { toc } { \l_tmpb_tl }
	\skip_horizontal:n { \l__zhvt_font_size_dim * 4 }
	\l_tmpa_tl
	\par
	\thispagestyle { empty }
}
\cs_set_protected:Npn \__zhvt_set_shukou_content: {
	% ★ 鱼尾上方留白（title 静态框）嵌入闲叟钤印：
	%   整页 /Rotate 90，故此处把图用 \rotatebox{90} 转正为竖章；
	%   尺寸卡在版心条幅内（书内约 7mm 宽 × 33.8mm 高，留余量不压鱼尾）。
	\setstaticcontents * { title } {
		\vbox_to_ht:nn { \l__zhvt_baseline_skip_dim } {
			\vskip 14.2pt
			\hbox_to_wd:nn { \l__zhvt_shukou_title_height_dim } {
				\hfil
				\rotatebox [ origin = c ] { 90 } { \includegraphics [ width = 7mm , height = 33.8mm ] { yshs.png } }
				\hfil
			}
		}
	}
% ★ 版心（书口）统一改法（2026-11-16 愚大人定）：18pt＋框内居中＋书口文字同版框色
%   ① 原 27.2pt 时 7 字长篇名（如「廉颇蔺相如列传」）末字溢出 chapter 框、压住页码；
%      改 \fontsize{18pt}{18pt}，前后 \hfil 令字串在框内居中（不随篇名字数漂移）。
%   ② \vbox{\vskip 8.2pt} 定量下移，抵消缩字号后栏宽方向 8.2pt 上偏（像素实测 y 中心 543.3 = 原版）。
%      ⚠ \raisebox / \vspace* 在 flowfram 框内无效；换字号须同步换补偿：16pt→10.1、19pt→7.4。
%   ③ chapter 与 page~number 均置 \color{sealbrown}（HTML 874434），与版框／界行同色；正文仍黑。
%   ⚠ 页码字位、钤印、鱼尾、界行、正文一律未动；五部合订本同此口径。
	\setdynamiccontents * { chapter }
	{ \color{sealbrown} \vbox { \vskip 8.2pt \hbox to \l__zhvt_shukou_chapter_title_height_dim
		{ \hfil \fontsize{18pt}{18pt}\selectfont \l__zhvt_title_str \quad \l__zhvt_chapter_title_tl \hfil } } }
	\setdynamiccontents * { page~number } { \color{sealbrown} \quad \zhnum { page } }
}
\__zhvt_set_shukou_content:
\ExplSyntaxOff

\maketitle

% 单篇，无多卷目录，不渲染目录页

\juan{日者列传}{卷一百二十七~日者列传第六十七}
\hspace*{2em}自古受命而王。王者之兴何尝不以卜筮决於天命哉。其於周尤甚。及秦可见。代王之入。任於卜者。太卜之起。由汉兴而有。\\
\hspace*{2em}司马季主者。楚人也。卜於长安东市。\\
\hspace*{2em}宋忠为中大夫。贾谊为博士。同日俱出洗沐。相从论议。诵易先王圣人之道术。究遍人情。相视而叹。贾谊曰。﹁吾闻古之圣人。不居朝廷。必在卜医之中。今吾已见三公九卿朝士大夫。皆可知矣。试之卜数中以观采。﹂二人即同舆而之市。游於卜肆中。天新雨。道少人。司马季主间坐。弟子三四人侍。方辩天地之道。日月之运。阴阳吉凶之本。二大夫再拜谒。司马季主视其状貌。如类有知者。即礼之。使弟子延之坐。坐定。司马季主复理前语。分别天地之终始。日月星辰之纪。差次仁义之际。列吉凶之符。语数千言。莫不顺理。\\
\hspace*{2em}宋忠。贾谊瞿然而悟。猎缨正襟危坐。曰。﹁吾望先生之状。听先生之辞。小子窃观於世。未尝见也。今何居之卑。何行之汙。﹂\\
\hspace*{2em}司马季主捧腹大笑曰。﹁观大夫类有道术者。今何言之陋也。何辞之野也。今夫子所贤者何也。所高者谁也。今何以卑汙长者。﹂\\
\hspace*{2em}二君曰。﹁尊官厚禄。世之所高也。贤才处之。今所处非其地。故谓之卑。言不信。行不验。取不当。故谓之汙。夫卜筮者。世俗之所贱简也。世皆言曰。﹃夫卜者多言夸严以得人情。虚高人禄命以说人志。擅言祸灾以伤人心。矫言鬼神以尽人财。厚求拜谢以私於己。﹄此吾之所耻。故谓之卑汙也。﹂\\
\hspace*{2em}司马季主曰。公且安坐。公见夫被发童子乎。日月照之则行。不照则止。问之日月疵瑕吉凶。则不能理。由是观之。能知别贤与不肖者寡矣。\\
\hspace*{2em}贤之行也。直道以正谏。三谏不听则退。其誉人也不望其报。恶人也不顾其怨。以便国家利众为务。故官非其任不处也。禄非其功不受也。见人不正。虽贵不敬也。见人有污。虽尊不下也。得不为喜。去不为恨。非其罪也。虽累辱而不愧也。\\
\hspace*{2em}今公所谓贤者。皆可为羞矣。卑疵而前。孅韱趋而言。相引以势。相导以利。比周宾正。以求尊誉。以受公奉。事私利。枉主法。猎农民。以官为威。以法为机。求利逆暴。譬无异於操白刃劫人者也。初试官时。倍力为巧诈。饰虚功执空文以诳主上。用居上为右。试官不让贤陈功。见伪增实。以无为有。以少为多。以求便势尊位。食饮驱驰。从姬歌兒。不顾於亲。犯法害民。虚公家。此夫为盗不操矛弧者也。攻而不用弦刃者也。欺父母未有罪而弑君未伐者也。何以为高贤才乎。\\
\hspace*{2em}盗贼发不能禁。夷貊不服不能摄。奸邪起不能塞。官秏乱不能治。四时不和不能调。岁穀不孰不能適。才贤不为。是不忠也。才不贤而讬官位。利上奉。妨贤者处。是窃位也。有人者进。有财者礼。是伪也。子独不见鸱枭之与凤皇翔乎。兰芷芎藭弃於广野。蒿萧成林。使君子退而不显众。公等是也。\\
\hspace*{2em}述而不作。君子义也。今夫卜者。必法天地。象四时。顺於仁义。分策定卦。旋式正釭。然後言天地之利害。事之成败。昔先王之定国家。必先龟策日月。而後乃敢代。正时日。乃後入家。产子必先占吉凶。後乃有之。自伏羲作八卦。周文王演三百八十四爻而天下治。越王句践放文王八卦以破敌国。霸天下。由是言之。卜筮有何负哉。\\
\hspace*{2em}且夫卜筮者。埽除设坐。正其冠带。然後乃言事。此有礼也。言而鬼神或以飨。忠臣以事其上。孝子以养其亲。慈父以畜其子。此有德者也。而以义置数十百钱。病者或以愈。且死或以生。患或以免。事或以成。嫁子娶妇或以养生。此之为德。岂直数十百钱哉。此夫老子所谓﹃上德不德。是以有德﹄。今夫卜筮者利大而谢少。老子之云岂异於是乎。\\
\hspace*{2em}庄子曰。﹃君子内无饥寒之患。外无劫夺之忧。居上而敬。居下不为害。君子之道也。﹄今夫卜筮者之为业也。积之无委聚。藏之不用府库。徙之不用辎车。负装之不重。止而用之无尽索之时。持不尽索之物。游於无穷之世。虽庄氏之行未能增於是也。子何故而云不可卜哉。天不足西北。星辰西北移。地不足东南。以海为池。日中必移。月满必亏。先王之道。乍存乍亡。公责卜者言必信。不亦惑乎。\\
\hspace*{2em}公见夫谈士辩人乎。虑事定计。必是人也。然不能以一言说人主意。故言必称先王。语必道上古。虑事定计。饰先王之成功。语其败害。以恐喜人主之志。以求其欲。多言夸严。莫大於此矣。然欲彊国成功。尽忠於上。非此不立。今夫卜者。导惑教愚也。夫愚惑之人。岂能以一言而知之哉。言不厌多。\\
\hspace*{2em}﹁故骐骥不能与罢驴为驷。而凤皇不与燕雀为群。而贤者亦不与不肖者同列。故君子处卑隐以辟众。自匿以辟伦。微见德顺以除群害。以明天性。助上养下。多其功利。不求尊誉。公之等喁喁者也。何知长者之道乎。﹂\\
\hspace*{2em}宋忠。贾谊忽而自失。芒乎无色。怅然噤口不能言。於是摄衣而起。再拜而辞。行洋洋也。出门仅能自上车。伏轼低头。卒不能出气。\\
\hspace*{2em}居三日。宋忠见贾谊於殿门外。乃相引屏语相谓自叹曰。﹁道高益安。势高益危。居赫赫之势。失身且有日矣。夫卜而有不审。不见夺糈。为人主计而不审。身无所处。此相去远矣。犹天冠地屦也。此老子之所谓﹃无名者万物之始﹄也。天地旷旷。物之熙熙。或安或危。莫知居之。我与若。何足预彼哉。彼久而愈安。虽曾氏之义未有以异也。﹂\\
\hspace*{2em}久之。宋忠使匈奴。不至而还。抵罪。而贾谊为梁怀王傅。王堕马薨。谊不食。毒恨而死。此务华绝根者也。\\
\hspace*{2em}太史公曰。古者卜人所以不载者。多不见于篇。及至司马季主。余志而著之。\\
\hspace*{2em}褚先生曰。臣为郎时。游观长安中。见卜筮之贤大夫。观其起居行步。坐起自动。誓正其衣冠而当乡人也。有君子之风。见性好解妇来卜。对之颜色严振。未尝见齿而笑也。从古以来。贤者避世。有居止舞泽者。有居民间闭口不言。有隐居卜筮间以全身者。夫司马季主者。楚贤大夫。游学长安。通易经。术黄帝。老子。博闻远见。观其对二大夫贵人之谈言。称引古明王圣人道。固非浅闻小数之能。及卜筮立名声千里者。各往往而在。传曰。﹁富为上。贵次之。既贵各各学一伎能立其身。﹂黄直。大夫也。陈君夫。妇人也。以相马立名天下。齐张仲。曲成侯以善击刺学用剑。立名天下。留长孺以相彘立名。荥阳褚氏以相牛立名。能以伎能立名者甚多。皆有高世绝人之风。何可胜言。故曰。﹁非其地。树之不生。非其意。教之不成。﹂夫家之教子孙。当视其所以好。好含苟生活之道。因而成之。故曰。﹁制宅命子。足以观士。子有处所。可谓贤人。﹂\\
\hspace*{2em}臣为郎时。与太卜待诏为郎者同署。言曰。﹁孝武帝时。聚会占家问之。某日可取妇乎。五行家曰可。堪舆家曰不可。建除家曰不吉。丛辰家曰大凶。历家曰小凶。天人家曰小吉。太一家曰大吉。辩讼不决。以状闻。制曰。﹃避诸死忌。以五行为主。﹄﹂人取於五行者也。\\

% ===== 索隐述赞（司马贞《史记索隐》）· 统一样式 =====
% 题头「索隐述赞曰」五字并入反白框；正文每四句一列，定宽盒 \hfil 居列中；
% \fontsize{24}{41} 列距对齐全书界行（121+41k）；\clearpage 使述赞独占一页。
\clearpage
\shuzanhead{索隐述赞曰}\linebreak
{\shuzanbody\fontsize{24}{41}\selectfont
\hbox to \hsize{\hfil 日者之名。有自来矣。吉凶占候。著於墨子。\hfil}\linebreak
\hbox to \hsize{\hfil 齐楚异法。书亡罕纪。後人斯继。季主独美。\hfil}\linebreak
\hbox to \hsize{\hfil 取免暴秦。此焉终否。\hfil}%
\par}
\end{document}
